% A LaTeX template for MSc Thesis submissions to 
% Politecnico di Milano (PoliMi) - School of Industrial and Information Engineering
%
% S. Bonetti, A. Gruttadauria, G. Mescolini, A. Zingaro
% e-mail: template-tesi-ingind@polimi.it
%
% Last Revision: October 2021
%
% Copyright 2021 Politecnico di Milano, Italy. NC-BY

\documentclass{Configuration_Files/PoliMi3i_thesis}

%------------------------------------------------------------------------------
%	REQUIRED PACKAGES AND  CONFIGURATIONS
%------------------------------------------------------------------------------

% CONFIGURATIONS
\usepackage{parskip} % For paragraph layout
\usepackage{setspace} % For using single or double spacing
\usepackage{emptypage} % To insert empty pages
\usepackage{multicol} % To write in multiple columns (executive summary)
\setlength\columnsep{15pt} % Column separation in executive summary
\setlength\parindent{0pt} % Indentation
\raggedbottom  

% PACKAGES FOR TITLES
\usepackage{titlesec}
% \titlespacing{\section}{left spacing}{before spacing}{after spacing}
\titlespacing{\section}{0pt}{3.3ex}{2ex}
\titlespacing{\subsection}{0pt}{3.3ex}{1.65ex}
\titlespacing{\subsubsection}{0pt}{3.3ex}{1ex}
\usepackage{color}

% PACKAGES FOR LANGUAGE AND FONT
\usepackage[english]{babel} % The document is in English  
\usepackage[utf8]{inputenc} % UTF8 encoding
\usepackage[T1]{fontenc} % Font encoding
\usepackage[11pt]{moresize} % Big fonts

% PACKAGES FOR IMAGES
\usepackage{graphicx}
\usepackage{transparent} % Enables transparent images
\usepackage{eso-pic} % For the background picture on the title page
\usepackage{subfig} % Numbered and caption subfigures using \subfloat.
\usepackage{tikz} % A package for high-quality hand-made figures.
\usetikzlibrary{}
\graphicspath{{./Images/}} % Directory of the images
\usepackage{caption} % Coloured captions
\usepackage{xcolor} % Coloured captions
\usepackage{amsthm,thmtools,xcolor} % Coloured "Theorem"
\usepackage{float}

% STANDARD MATH PACKAGES
\usepackage{amsmath}
\usepackage{amsthm}
\usepackage{amssymb}
\usepackage{amsfonts}
\usepackage{bm}
\usepackage[overload]{empheq} % For braced-style systems of equations.
\usepackage{fix-cm} % To override original LaTeX restrictions on sizes

% PACKAGES FOR TABLES
\usepackage{tabularx}
\usepackage{longtable} % Tables that can span several pages
\usepackage{colortbl}
\usepackage{longtable}
\usepackage{geometry}

% PACKAGES FOR ALGORITHMS (PSEUDO-CODE)
\usepackage{algorithm}
\usepackage{algorithmic}

% PACKAGES FOR REFERENCES & BIBLIOGRAPHY
\usepackage[colorlinks=true,linkcolor=black,anchorcolor=black,citecolor=black,filecolor=black,menucolor=black,runcolor=black,urlcolor=black]{hyperref} % Adds clickable links at references
\usepackage{cleveref}
\usepackage[square, numbers, sort&compress]{natbib} % Square brackets, citing references with numbers, citations sorted by appearance in the text and compressed
\bibliographystyle{abbrvnat} % You may use a different style adapted to your field

% OTHER PACKAGES
\usepackage{pdfpages} % To include a pdf file
\usepackage{afterpage}
\usepackage{lipsum} % DUMMY PACKAGE
\usepackage{fancyhdr} % For the headers
\fancyhf{}

% Input of configuration file. Do not change config.tex file unless you really know what you are doing. 
% Define blue color typical of polimi
\definecolor{bluepoli}{cmyk}{0.4,0.1,0,0.4}

% Custom theorem environments
\declaretheoremstyle[
  headfont=\color{bluepoli}\normalfont\bfseries,
  bodyfont=\color{black}\normalfont\itshape,
]{colored}

% Set-up caption colors
\captionsetup[figure]{labelfont={color=bluepoli}} % Set colour of the captions
\captionsetup[table]{labelfont={color=bluepoli}} % Set colour of the captions
\captionsetup[algorithm]{labelfont={color=bluepoli}} % Set colour of the captions

\theoremstyle{colored}
\newtheorem{theorem}{Theorem}[chapter]
\newtheorem{proposition}{Proposition}[chapter]

% Enhances the features of the standard "table" and "tabular" environments.
\newcommand\T{\rule{0pt}{2.6ex}}
\newcommand\B{\rule[-1.2ex]{0pt}{0pt}}

% Pseudo-code algorithm descriptions.
\newcounter{algsubstate}
\renewcommand{\thealgsubstate}{\alph{algsubstate}}
\newenvironment{algsubstates}
  {\setcounter{algsubstate}{0}%
   \renewcommand{\STATE}{%
     \stepcounter{algsubstate}%
     \Statex {\small\thealgsubstate:}\space}}
  {}

% New font size
\newcommand\numfontsize{\@setfontsize\Huge{200}{60}}

% Title format: chapter
\titleformat{\chapter}[hang]{
\fontsize{50}{20}\selectfont\bfseries\filright}{\textcolor{bluepoli} \thechapter\hsp\hspace{2mm}\textcolor{bluepoli}{|   }\hsp}{0pt}{\huge\bfseries \textcolor{bluepoli}
}

% Title format: section
\titleformat{\section}
{\color{bluepoli}\normalfont\Large\bfseries}
{\color{bluepoli}\thesection.}{1em}{}

% Title format: subsection
\titleformat{\subsection}
{\color{bluepoli}\normalfont\large\bfseries}
{\color{bluepoli}\thesubsection.}{1em}{}

% Title format: subsubsection
\titleformat{\subsubsection}
{\color{bluepoli}\normalfont\large\bfseries}
{\color{bluepoli}\thesubsubsection.}{1em}{}

% Shortening for setting no horizontal-spacing
\newcommand{\hsp}{\hspace{0pt}}

\makeatletter
% Renewcommand: cleardoublepage including the background pic
\renewcommand*\cleardoublepage{%
  \clearpage\if@twoside\ifodd\c@page\else
  \null
  \AddToShipoutPicture*{\BackgroundPic}
  \thispagestyle{empty}%
  \newpage
  \if@twocolumn\hbox{}\newpage\fi\fi\fi}
\makeatother

%For correctly numbering algorithms
\numberwithin{algorithm}{chapter}

%----------------------------------------------------------------------------
%	NEW COMMANDS DEFINED
%----------------------------------------------------------------------------

% EXAMPLES OF NEW COMMANDS
\newcommand{\bea}{\begin{eqnarray}} % Shortcut for equation arrays
\newcommand{\eea}{\end{eqnarray}}
\newcommand{\e}[1]{\times 10^{#1}}  % Powers of 10 notation

%----------------------------------------------------------------------------
%	ADD YOUR PACKAGES (be careful of package interaction)
%----------------------------------------------------------------------------

%----------------------------------------------------------------------------
%	ADD YOUR DEFINITIONS AND COMMANDS (be careful of existing commands)
%----------------------------------------------------------------------------

%----------------------------------------------------------------------------
%	BEGIN OF YOUR DOCUMENT
%----------------------------------------------------------------------------

\begin{document}

\fancypagestyle{plain}{%
\fancyhf{} % Clear all header and footer fields
\fancyhead[RO,RE]{\thepage} %RO=right odd, RE=right even
\renewcommand{\headrulewidth}{0pt}
\renewcommand{\footrulewidth}{0pt}}

%----------------------------------------------------------------------------
%	TITLE PAGE
%----------------------------------------------------------------------------

\pagestyle{empty} % No page numbers
\frontmatter % Use roman page numbering style (i, ii, iii, iv...) for the preamble pages

\puttitle{
	title=Software Engineering 2 - RASD, % Title of the thesis
	name=Simone Controguerra,
        name=Matteo Delogu,
            name=Gianluca Di Paola,% Author Name and Surname % Student ID number (numero di matricola)
	advisor= Prof. Elisabetta Di Nitto, % Supervisor name
	 % Co-Supervisor name, remove this line if there is none
	academicyear={2025-2026}  % Academic Year
} 

%----------------------------------------------------------------------------
%	PREAMBLE PAGES: ABSTRACT (inglese e italiano), EXECUTIVE SUMMARY
%----------------------------------------------------------------------------
\startpreamble
\setcounter{page}{1} % Set page counter to 1

%----------------------------------------------------------------------------
%	LIST OF CONTENTS/FIGURES/TABLES/SYMBOLS
%----------------------------------------------------------------------------

% TABLE OF CONTENTS
\thispagestyle{empty}
\tableofcontents % Table of contents 
\thispagestyle{empty}
\cleardoublepage

%-------------------------------------------------------------------------
%	THESIS MAIN TEXT
%-------------------------------------------------------------------------
% In the main text of your thesis you can write the chapters in two different ways:
%
%(1) As presented in this template you can write:
%    \chapter{Title of the chapter}
%    *body of the chapter*
%
%(2) You can write your chapter in a separated .tex file and then include it in the main file with the following command:
%    \chapter{Title of the chapter}
%    \input{chapter_file.tex}
%
% Especially for long thesis, we recommend you the second option.

\addtocontents{toc}{\vspace{2em}} % Add a gap in the Contents, for aesthetics
\mainmatter % Begin numeric (1,2,3...) page numbering

% --------------------------------------------------------------------------
% NUMBERED CHAPTERS % Regular chapters following
% --------------------------------------------------------------------------

\chapter{Introduction}
\label{ch:chapter_one}%
% The \label{...}% enables to remove the small indentation that is generated, always leave the % symbol.



\section{Purpose}


\section{Scope}


\subsection{Goals}
\begin{itemize}
    \item \textbf{[G1]} Enable the cycling community to access a shared catalog of safe and reliable bike paths.
	\item \textbf{[G2]} Allow registered users to store trips with automatically-collected sensor data and associated weather (if available).
	\item \textbf{[G3]} Develop and maintain a collaborative inventory of bike paths and their score, along with status and obstacle information regarding the streets they are composed of.
	\item \textbf{[G4]} Enable registered users to contribute to the enrichment of the bike paths catalog, by publishing automatically-recorded bike paths.
	\item \textbf{[G5]} Enable registered users to contribute to the enrichment of the bike paths catalog, by publishing manually-entered bike paths.
	\item \textbf{[G6]} Allow any user (registered or unregistered) to search for bike paths from Origin to Destination, ordered by their effectiveness for such parameters.
	\item \textbf{[G7]} Ensure bike paths consist only of roads suitable and safe for cycling.
	\item \textbf{[G8]} Reduce errors and false positives through user confirmation and information merging.
\end{itemize}

\subsection{Phenomena Table}


\begin{center}
    \begin{longtable}{|c|p{7cm}|c|c|}
    
        \hline
        \textbf{Code} & \textbf{Phenomenon} & \textbf{Who controls it?} & \textbf{Shared?} \\ \hline
        \endfirsthead
        
        \hline
        \textbf{Code} & \textbf{Phenomenon} & \textbf{Who controls it} & \textbf{Shared?} \\ \hline
        \endhead
        
        \hline
        % **WORLD ONLY (not shared)**
         \textbf{[P1]} & Any user (registered or not) follow a new bike path. & World & No\\ \hline
         \textbf{[P2]} & A registered user starts cycling. & World & No \\ \hline
         \textbf{[P3]} & Road status (needs maintenence, optimal, mediocre, ...). & World & No \\ \hline
        
        % **SHARED, WORLD-CONTROLLED**
         \textbf{[P4]} & A new user creates a profile. & World & Yes \\ \hline
         \textbf{[P5]} & A registered user starts a trip session. & World & Yes \\ \hline
         \textbf{[P6]} & A registered user manually defines (in terms of streets, status, etc.) and publishes a bike path. & World & Yes \\ \hline
         \textbf{[P7]} & A registered user confirms or corrects information acquired by the system and publishes the autogenerated bike path. & World & Yes \\ \hline
         \textbf{[P8]} & Any user (registered or not) requires a bike path by specifying origin and destination. & World & Yes \\ \hline
         \textbf{[P9]} & Device sensors (GPS, accelerometer, gyroscope) acquire data. & World & Yes \\ \hline
        \textbf{[P10]} & The weather-service responds to the request. & World & Yes \\ \hline
        \textbf{[P11]} & The navigation-system service responds to the requests. & World & Yes \\ \hline
        
        % **SHARED, MACHINE-CONTROLLED**
        \textbf{[P12]} & The system visually notifies the user that the trip recording has started. & Machine & Yes \\ \hline
        \textbf{[P13]} & The system provides all recorded trip metrics. & Machine & Yes \\ \hline
        \textbf{[P14]} & The system notifies the registered user that an autogenerated bike path is ready for detection revision. & Machine & Yes \\ \hline
        \textbf{[P15]} & The system shows to any user (registered or not) all the possible bike paths, ordered by overall score, to accomplish their request. & Machine & Yes \\ \hline
        \textbf{[P16]} & The system prevents registered users from publishing bike paths with unallowed streets. & Machine & Yes \\ \hline
        \textbf{[P17]} & The system updates the score of a bike path whenever the status of a street changes. & Machine & Yes \\ \hline
        \textbf{[P18]} & The system calls the weather-forecast service to add weather info to a trip. & Machine & Yes \\ \hline
        \textbf{[P19]} & The system calls the navigation-system service to evaluate road safety for bikers (elements taken into account are of the likings of: proper bike track, car presence frequence, average speed of a bike, etc.). & Machine & Yes \\ \hline
        
        % **MACHINE ONLY (not shared)**
        \textbf{[P20]} & The system detects when the registered user is cycling and processes relevant data, to obtain performance metrics, to detect obstacles and such. & Machine & No \\ \hline
        \textbf{[P21]} & The system calculates the overall score of a bike path, depending on its road status and effectiveness (in terms of allowing any user to reach the desired destination from a given origin). & Machine & No \\ \hline
        \textbf{[P22]} & The system re-computes the score of a bike path whenever the status of a street changes. & Machine & No \\ \hline
        \textbf{[P23]} & The system stores a new bike path. & Machine & No \\ \hline
        
    \caption{Phenomena}
    \end{longtable}
\end{center}

\vfill

\section{Definitions, acronyms and abbreviations}

\begin{itemize}

    \item \textbf{API (Application Programming Interface):} A set of protocols and tools that enable different software applications to communicate with each other.
    \item \textbf{Redundancy:} The inclusion of supplementary components or systems to ensure reliability and availability in the event of a primary component failure.
    \item \textbf{Scalability:} The system's capacity to manage a growing number of users or increased load by incorporating additional resources without significant degradation in performance.
    \item \textbf{UI (User Interface):} The interface through which users interact with the platform, encompassing buttons, screens, and navigation elements.
    \item \textbf{Uptime:} The period during which a system is fully operational and accessible to users without interruption.
    \item \textbf{Backup:} The process of duplicating and archiving data to ensure restorability in the event of loss or damage to the original data.
    \item \textbf{JSON (JavaScript Object Notation):} A lightweight data-interchange format that is easily readable and writable by both humans and machines.
    \item \textbf{Bike Path:} A route specifically designated or suitable for bicycle travel, which can be either manually defined by users or automatically generated by the system. It is formally defined as a set consisting of: a Starting Point, a Finishing Point, and an ordered sequence of traversable Streets (see definition of Street).
    \item \textbf{Trip:} A standalone, recorded instance of a user traveling from one point to another, which includes sensor data, metrics, and potentially weather information.
    \item \textbf{Street:} A segment of road suitable for cycling (factors taken into account include average bike speed limit, roadway width, traffic volume, dedicated bike lanes, etc.) that forms part of a bike path, with attributes such as status and presence of irregularities (e.g., potholes).
    \item \textbf{Navigation System:} An external service or API used by the system to evaluate road safety for bikers, considering factors such as bike track availability and traffic conditions.
    \item \textbf{Weather-forecast System:} An external service or API used by the system to provide weather information associated with a trip.

\end{itemize}

\section{Revision history}


\section{Reference documents}
Professor Elisabetta Di Nitto, \textit{Slides of the course "Software Engineering 2"} at PoliMi.
\\ 
Polytechnic University of Milan, \textit{Specifications of the RASD assignment for the Software Engineering II course}.

\section{Document structure}



\thispagestyle{empty}
\cleardoublepage




%------------------------------------------------------

\chapter{Overall Description}


\section{Product perspective}


\subsection{Scenarios}

\subsection{Domain class diagram}


\section{Product functions}


\section{User characteristics}


\section{Assumptions, dependencies and constraints}

In order to satisfy all the goals, the system requires the following assumptions to be met:
\begin{itemize}  
	\item \textbf{[A1] Compliance with guidelines:} All users (registered and unregistered) are expected to follow the platform’s rules and guidelines and provide accurate, up-to-date information.
	\item \textbf{[A2] Valid credentials:} To use the service, registered users are expected to have a valid email address.
	\item \textbf{[A3] User’s mobile device:} The user must possess a mobile device (e.g., smartphone) with integrated sensors (mandatory: GPS, accelerometer, gyroscope) capable of internet communication and application processing.
	\item \textbf{[A4] Sensor reliability:} GPS, accelerometer, gyroscope, and other necessary sensors are available, functioning correctly, and calibrated on the user's device.
	\item \textbf{[A5] GPS data detection:} The error margin of the GPS signal shall be within acceptable thresholds to allow for the detection of the starting point (origin) and arrival point (destination).
	\item \textbf{[A6] Accelerometer and Gyroscope sensitivity:} Data detection from accelerometer and gyroscope sensors must be accurate enough to correctly detect potholes and other street irregularities.
\end{itemize}
In addition, the systems requires some certain dependencies:
\begin{itemize}  
	\item \textbf{[D1] Reliable internet access:} All users (registered and unregistered) shall have a stable internet connection to perform certain actions.
	\item \textbf{[D2] Weather-forecast Service API:} The weather forecast service API must be correctly implemented, sufficiently available, and functioning correctly.
	\item \textbf{[D3] Navigation-system Service API:} The navigation system service API must be correctly implemented, sufficiently available, and functioning correctly.
\end{itemize}

%------------------------------------------------------
\chapter{Specific Requirements}

\section{External Interface requirements}

\subsection{User interfaces}

\subsection{Hardware Interfaces}

\subsection{Software interfaces}

\subsection{Communication interfaces}

\section{Functional Requirements}

\subsection{Requirements Mapping}

\subsection{Use case diagrams}

\subsection{Use Cases}

\section{Performance Requirements}

\section{Design Constraints}

\subsection{Standards compliance}

\subsection{Hardware limitations}

\subsection{Any other constraint }

\section{Software system attributes}

\subsection{Reliability}

\subsection{Availability}

\subsection{Security}

\subsection{Maintainability}

\subsection{Portability }

%------------------------------------------------------
\chapter{Formal Analysis using Alloy}


\section{Examples}

%------------------------------------------------------
\chapter{Effort spent}

\chapter{References}


%------------------------------------------------------
\chapter*{List of Figures}


\chapter*{List of Tables}


%------------------------------------------------------
%------------------------------------------------------
\cleardoublepage
\end{document}
